\documentclass{article}
\usepackage{amsmath,amssymb,xurl}
\usepackage[hidelinks]{hyperref}
% Shared commands for notes in docs/paper-gaps/.

\DeclareUnicodeCharacter{2080}{\ifmmode{}_{0}\else\textsubscript{0}\fi}
\DeclareUnicodeCharacter{2081}{\ifmmode{}_{1}\else\textsubscript{1}\fi}
\DeclareUnicodeCharacter{2082}{\ifmmode{}_{2}\else\textsubscript{2}\fi}
\DeclareUnicodeCharacter{2099}{\ifmmode{}_{n}\else\textsubscript{n}\fi}

\newcommand{\Lean}{\textsc{Lean}}
\ExplSyntaxOn
\NewDocumentCommand{\leanid}{m}{
  \ifmmode
    \textsf{\detokenize{#1}}
  \else
    \group_begin:
    \urlstyle{sf}
    \tl_set:Nn \l_tmpa_tl {#1}
    \tl_replace_all:Nnn \l_tmpa_tl {\_} {_}
    \exp_args:NV \path \l_tmpa_tl
    \group_end:
  \fi
}
\ExplSyntaxOff
\providecommand{\tr}{\operatorname{tr}}
\newcommand{\tnleanIssueBase}{https://github.com/LionSR/TNLean/issues/}
\newcommand{\tnleanPrBase}{https://github.com/LionSR/TNLean/pull/}
\newcommand{\tnleanDocBase}{https://sirui-lu.com/TNLean/docs/}
\newcommand{\ghissue}[1]{\href{\tnleanIssueBase#1}{GitHub issue \##1}}
\newcommand{\ghpr}[1]{\href{\tnleanPrBase#1}{PR \##1}}
\newcommand{\doclink}[1]{\href{\tnleanDocBase#1.html}{\path{#1}}}

% Dirac notation and the identity operator, as in blueprint/src/macros/common.tex.
% \providecommand keeps note-local definitions authoritative.
\usepackage{dsfont}
\providecommand{\ket}[1]{|#1\rangle}
\providecommand{\bra}[1]{\langle #1|}
\providecommand{\braket}[2]{\langle #1 | #2 \rangle}
\providecommand{\ketbra}[2]{|#1\rangle\!\langle #2|}
\providecommand{\Id}{\mathds{1}}
\providecommand{\one}{\mathds{1}}
\providecommand{\C}{\mathbb{C}}
\providecommand{\MN}[1]{M_{#1}(\C)}
\providecommand{\MD}{M_D(\C)}

% Verdict marker: \gapnote{<kind>}{<status>}. Kind is the policy
% classification (clarification, local-correction, scope-restriction,
% unfaithful, false-source, open-gap); status is open, wip (elimination
% underway), resolved, or historical. Machine-read by the site index;
% prints a verdict line.
\newcommand{\gapnote}[2]{\par\noindent\textbf{Verdict:} #1 (#2).\par}

\title{A Finite Zero-Crossing Criterion for Native Flux Vacuum Density}
\author{}
\date{}
\begin{document}
\maketitle
\gapnote{scope-restriction}{open}

For the local-density comparison in Schuch, Cirac, and P\'erez-Garc\'ia,
\href{https://arxiv.org/abs/1001.3807v3}{arXiv:1001.3807v3}, Theorem 6.9
and Lemma 6.14, source \path{Papers/1001.3807/paper_v3.tex}, lines
1935--1990 and 2199--2214, one can use conditions on the actual finite
crossing graph in place of a simply connected envelope.

The graph \leanid{TNLean.PEPS.torusDualZeroCrossingGraph} retains exactly
the native physical bonds with zero signed crossing count for the supplied
actual path. Both edge orientations are included and the periodic seams
are treated in the native coordinates. The graph is independent of the
insertion group and element. Its edges have identity actual flux transport.

The theorem
\leanid{TNLean.PEPS.IsGIsometric.torusFlux_normalized_density_eq_vacuum_of_zero_crossing_connectivity}
assumes connected observation and complementary induced graphs, nonempty
boundary, zero signed counts on observed incident bonds, and one component
of this computed zero-crossing graph in the complement containing all
complementary boundary endpoints. It derives both actual cuts nonzero and
normalized observed density equal to vacuum for every insertion element
of a finite group and every native regular group-isometric tensor.
Both periods are at least three. No simple connectedness, envelope,
supplied untwisted subgraph, gauge, factorization, Gram identity, density
equality or nonzero witness is assumed.

These are finite graph and integer-count conditions on the actual path.
They are sufficient, not necessary for density equality, and no terminating
classification algorithm for all observations is asserted. In particular,
the observed and complementary graphs must still be connected; an annulus
with disconnected complement fails that hypothesis. The positive native
simple-connectedness theorem and its envelope extension remain unchanged.
The explicit annular obstruction prevents an endpoint-only generalization.
The broader classification remains open under \ghissue{8703}; no
\ghissue{8947} Hamiltonian files or claims are changed.
\end{document}
